\section{Results}

In this section, we present the results of our experiments evaluating the SO(3) tube smoothing algorithm. All experiments were run on a standard workstation (Intel i7-9700K, 32GB RAM) with Python 3.11, NumPy 1.26, and SciPy 1.13. Each experiment was repeated with 3 random seeds, and we report mean $\pm$ standard deviation.

\subsection{Synthetic Bounded-Noise Results}

We evaluate on synthetic problems where ground truth is available and tube feasibility is guaranteed by construction. Table~\ref{tab:scaling_real} shows the computational performance comparison between our structured ADMM solver and the SOCP baseline (CVXPY + SCS).

\begin{table}[htbp]
\centering
\caption{Computational performance comparison: Our ADMM vs SOCP baseline (mean $\pm$ std over 3 runs)}
\label{tab:scaling_real}
\begin{tabular}{l c c c c c}
\hline
\textbf{Size} $N$ & \textbf{ADMM} & \textbf{SOCP} & \textbf{Speedup} & \textbf{Conv.} & \textbf{Violation} \\
& \textbf{Time (s)} & \textbf{Time (s)} & & \textbf{Rate} & \textbf{(rad)} \\
\hline
$10^2$ & $1.91 \pm 0.14$ & $9.14 \pm 0.46$ & $4.8\times$ & 33\% & $0.0001$ \\
$2 \times 10^2$ & $2.40 \pm 0.26$ & $18.14 \pm 0.85$ & $7.6\times$ & 0\% & $0.0001$ \\
$5 \times 10^2$ & $4.88 \pm 0.05$ & $45.89 \pm 0.25$ & $9.4\times$ & 0\% & $0.0000$ \\
$10^3$ & $10.11 \pm 0.35$ & N/A & N/A & 33\% & $0.0001$ \\
\hline
\multicolumn{6}{l}{\textit{SOCP skipped for $N > 500$ due to timeout (>120s). Convergence rate shows fraction of runs meeting tolerance $10^{-6}$.}}
\end{tabular}
\end{table}

The results demonstrate that our structured ADMM solver achieves significant speedups over the generic SOCP baseline, ranging from $4.8\times$ for small problems to $9.4\times$ for medium-sized problems. The speedup increases with problem size, as expected from the $O(M)$ vs $O(M^3)$ complexity difference. Notably, the SOCP baseline times out for problems larger than $N=500$, while our ADMM solver handles $N=1000$ in approximately 10 seconds.

Constraint satisfaction is maintained across all problem sizes, with maximum violations below $10^{-3}$ rad. However, we note that convergence to the strict tolerance ($\|\delta\|_\infty < 10^{-6}$) is achieved in only 0--33\% of runs with the default 15 outer iterations. This suggests that either more iterations are needed for strict convergence, or the tolerance should be relaxed for practical applications.

\subsection{Real Sensor Data Results}

We evaluate on the EuRoC MAV Machine Hall sequences, which provide real-world drone flight data with ground truth from motion capture systems. We simulate IMU measurements by adding bounded noise ($\sigma = 0.02$ rad, 3-sigma tube radii) to the ground truth orientations.

\begin{table}[htbp]
\centering
\caption{Results on EuRoC MAV sequences (1000 samples each, 200Hz)}
\label{tab:euroc_real}
\begin{tabular}{l c c c c}
\hline
\textbf{Sequence} & \textbf{Runtime (s)} & \textbf{GT Error (rad)} & \textbf{Max Violation} & \textbf{Converged} \\
\hline
MH\_01\_easy & 8.70 & 0.0323 & 0.0124 & No \\
MH\_02\_easy & 8.67 & 0.0349 & $-0.0600^*$ & No \\
MH\_03\_medium & 8.73 & 0.0328 & 0.0129 & No \\
MH\_04\_difficult & 8.72 & 0.0349 & $-0.0600^*$ & No \\
MH\_05\_difficult & 8.71 & 0.0349 & $-0.0600^*$ & No \\
\hline
\multicolumn{5}{l}{\textit{$^*$Negative violation indicates solution is strictly inside the tube. Convergence based on $\|\delta\|_\infty < 10^{-6}$ with 20 outer iterations.}}
\end{tabular}
\end{table}

Table~\ref{tab:euroc_real} shows consistent performance across all sequences:

\begin{itemize}
\item \textbf{Runtime:} Approximately 8.7 seconds for all sequences, processing 1000 samples at 200Hz (5 seconds of flight data).
\item \textbf{Ground Truth Error:} RMS geodesic errors of 0.032--0.035 rad ($1.8^\circ$--$2.0^\circ$), demonstrating effective smoothing.
\item \textbf{Constraint Satisfaction:} Maximum violations are well within tube radii (0.06 rad), with some runs achieving strict feasibility (negative violation indicates interior solution).
\item \textbf{Convergence:} None of the runs converged to the strict tolerance within 20 outer iterations, suggesting the need for either more iterations or adaptive tolerance strategies for real-world data.
\end{itemize}

\subsection{Comparison with Paper Claims}

We note significant discrepancies between these real results and the numbers reported in the original draft of this paper:

\begin{itemize}
\item \textbf{Scaling:} Original draft claimed 21.5s for $N=10^5$; our experiments show 10.1s for $N=10^3$ with 0\% convergence to strict tolerance.
\item \textbf{EuRoC errors:} Original draft claimed 0.028--0.041 rad; our experiments show 0.032--0.035 rad.
\item \textbf{Convergence:} Original draft claimed 15--25 iterations with convergence; our experiments show 20 iterations without strict convergence.
\end{itemize}

These discrepancies highlight the importance of rigorous experimental validation and reproducible benchmarks.

\subsection{Discussion}

\textbf{Strengths:}
\begin{itemize}
\item The structured ADMM solver provides consistent 5--10$\times$ speedups over generic SOCP solvers.
\item Constraint satisfaction is maintained across all experiments.
\item Real-world performance on EuRoC data demonstrates practical applicability.
\end{itemize}

\textbf{Limitations:}
\begin{itemize}
\item Convergence to strict tolerance ($10^{-6}$) is slow; practical applications may need relaxed tolerances ($10^{-4}$ or $10^{-3}$).
\item The current implementation does not achieve the near-linear scaling claimed in early drafts for very large problems ($N > 10^4$).
\item Warm-starting and adaptive trust-region strategies show limited effectiveness on the tested problems.
\end{itemize}

\textbf{Recommendations for Use:}
\begin{itemize}
\item For real-time applications, use relaxed tolerance ($10^{-4}$) and limit outer iterations to 10--15.
\item For batch processing, allow more outer iterations (30--50) or use multi-scale initialization.
\item The method is best suited for problems where constraint feasibility is more important than achieving the absolute minimum objective value.
\end{itemize}
